\chapter{\IfLanguageName{dutch}{Stand van zaken}{State of the art}}%
\label{ch:stand-van-zaken}

% Tip: Begin elk hoofdstuk met een paragraaf inleiding die beschrijft hoe
% dit hoofdstuk past binnen het geheel van de bachelorproef. Geef in het
% bijzonder aan wat de link is met het vorige en volgende hoofdstuk.

% Pas na deze inleidende paragraaf komt de eerste sectiehoofding.
%
% Dit hoofdstuk bevat je literatuurstudie. De inhoud gaat verder op de inleiding, maar zal het onderwerp van de bachelorproef *diepgaand* uitspitten. De bedoeling is dat de lezer na lezing van dit hoofdstuk helemaal op de hoogte is van de huidige stand van zaken (state-of-the-art) in het onderzoeksdomein. Iemand die niet vertrouwd is met het onderwerp, weet nu voldoende om de rest van het verhaal te kunnen volgen, zonder dat die er nog andere informatie moet over opzoeken 
%
%Je verwijst bij elke bewering die je doet, vakterm die je introduceert, enz.\ naar je bronnen. In \LaTeX{} kan dat met het commando Als argument van het commando geef je de ``sleutel'' van een ``record'' in een bibliografische databank in het Bib\LaTeX{}-formaat (een tekstbestand). Als je expliciet naar de auteur verwijst in de zin (narratieve referentie), gebruik je \texttt{$\backslash${}textcite\{\}}. Soms is de auteursnaam niet expliciet een onderdeel van de zin, dan gebruik je \texttt{$\backslash${}autocite\{\}} (referentie tussen haakjes). Dit gebruik je bv.~bij een citaat, of om in het bijschrift van een overgenomen afbeelding, broncode, tabel, enz. te verwijzen naar de bron. In de volgende paragraaf een voorbeeld van elk.
%
%
%Let er ook op: het \texttt{cite}-commando voor de punt, dus binnen de zin. Je verwijst meteen naar een bron in de eerste zin die erop gebaseerd is, dus niet pas op het einde van een paragraaf.


\subsection{Inleiding}

In dit hoofdstuk wordt de huidige stand van zaken rond large language models voor codegeneratie en hun toepassing in Python gebaseerde data analyse en ETL taken besproken. Waar de inleiding en probleemstelling vooral de context, doelgroep en onderzoeksvraag van deze bachelorproef schetsen, gaat dit hoofdstuk dieper in op de bestaande wetenschappelijke inzichten over de betrouwbaarheid van LLM gegenereerde code, typische fouttypes, gebruikte trainingsdata en reeds voorgestelde evaluatiekaders. Op basis van deze literatuurstudie wordt duidelijk welke aspecten al goed onderzocht zijn en waar nog tekortkomingen bestaan, in het bijzonder met betrekking tot data engineering scenario's die relevant zijn voor dit onderzoek. Deze analyse vormt de inhoudelijke basis voor het volgende hoofdstuk over de methodologie, waarin de concrete onderzoeksopzet, de gekozen taken, de evaluatiemetrics en de experimentele aanpak worden uitgewerkt in functie van de geformuleerde onderzoeksvraag en deelvragen.


\subsection{Large language models voor codegeneratie}

Large language models (LLM's) worden steeds vaker ingezet als algemene code assistenten die op basis van natuurlijke taal programmeercode genereren in talen zoals Python, Java en JavaScript \autocite{Li2023}. In tegenstelling tot traditionele code synthesis technieken, die vaak sterk afhankelijk zijn van formele specificaties of domeinspecifieke talen, maken LLM's gebruik van grootschalige taalmodellen die zijn getraind op enorme hoeveelheden broncode en natuurlijke taal, waardoor zij in staat zijn om volledige functies of scripts te produceren op basis van korte tekstuele beschrijvingen. \textcite{Li2023} tonen aan dat moderne modellen in veel gevallen syntactisch correcte en semantisch plausibele code genereren voor diverse programmeertaken. Dit verklaart deels waarom studenten en ontwikkelaars deze systemen spontaan omarmen als vervanging of aanvulling op traditionele bronnen zoals documentatie en Q\&A platforms.

Ondanks deze indrukwekkende capaciteiten blijft de betrouwbaarheid van LLM gegenereerde code een belangrijk aandachtspunt \autocite{Li2023}. LLM's optimaliseren niet expliciet voor correctheid, maar voor waarschijnlijkheid gegeven de trainingsdata, waardoor zij ook codefragmenten kunnen voorstellen die syntactisch juist zijn maar logische fouten bevatten, niet robuust zijn voor randgevallen of gebruikmaken van verouderde API's \autocite{Zhang2023}. In de context van Python betekent dit bijvoorbeeld dat een model een oplossing genereert die correct werkt voor een eenvoudig voorbeeld, maar faalt bij grotere datasets, ontbrekende waarden of niet voorziene inputformaten. Deze spanning tussen plausibiliteit en betrouwbaarheid vormt de kern van de problematiek die in deze bachelorproef wordt onderzocht.


\subsection{Betrouwbaarheid en defecten in LLM gegenereerde code}

Verschillende recente studies hebben de foutbestendigheid van LLM gegenereerde code systematisch onderzocht \autocite{Li2023,Esfahani2024}. \textcite{Li2023} evalueren de robuustheid en betrouwbaarheid van codegeneratie door LLM's  gegenereerde oplossingen te onderwerpen aan variaties in promptformulering en aanvullende testcases. Zij concluderen dat veel modellen gevoelig zijn voor kleine wijzigingen in de instructies, wat kan leiden tot uiteenlopende oplossingskwaliteiten, zelfs wanneer de functionele vereisten ongewijzigd blijven \autocite{Li2023}. Dit wijst erop dat de output van een LLM niet alleen afhangt van de taak zelf, maar ook van de manier waarop studenten hun vraag formuleren, wat in een onderwijscontext een extra bron van variabiliteit is.

\textcite{Esfahani2024} voeren een defectanalyse uit op code die door verschillende LLM's is gegenereerd, en classificeren de gevonden problemen in categorieën zoals syntaxisfouten, logische fouten, API misbruik, foutieve randgevalafhandeling en performantieproblemen. Zij rapporteren dat een aanzienlijk deel van de geanalyseerde codefragmenten minstens één defect bevat en dat logische  semantische fouten vaker voorkomen dan louter syntactische \autocite{Esfahani2024}. Dit is bijzonder relevant voor minder ervaren programmeurs, zoals studenten, omdat semantische fouten doorgaans moeilijker te detecteren zijn op basis van een oppervlakkige visuele inspectie. Bovendien tonen hun resultaten aan dat sommige modellen geneigd zijn om verouderde of niet meer aanbevolen functies te gebruiken, wat een risico vormt voor de kwaliteit en onderhoudbaarheid van de gegenereerde oplossingen.

Een bijkomende bevinding is dat LLM's vaak een grote mate van zelfvertrouwen uitstralen in de vorm van overtuigende toelichtingen bij hun code, ook wanneer de onderliggende oplossing fouten bevat \autocite{Li2023,Esfahani2024}. Dit fenomeen wordt vaak beschreven als "hallucinatie" of "confident wrong" wat gebruikers kan misleiden om foutieve code te accepteren zonder verdere tests of documentatie te raadplegen.

\subsection{Prompt engineering bij codegeneratie}
\label{sec:prompt_engineering}

Prompt engineering speelt een belangrijke rol bij het gebruik van large language models voor codegeneratie, omdat kleine verschillen in de formulering van een prompt een duidelijke invloed kunnen hebben op de kwaliteit van de gegenereerde code. \textcite{Guo2024} tonen in hun studie over codegeneratie met ChatGPT aan dat verschillende prompten leiden tot merkbare verschillen in de bruikbaarheid en correctheid van de output. Hun resultaten tonen dat prompts die het model vanuit een specifieke rol of met een duidelijke taakomschrijving aansturen betere resultaten opleveren dan algemene, weinig gestructureerde instructies \autocite{Guo2024}. Voor de context van deze bachelorproef is dit bijzonder relevant, omdat studenten vaak spontane en ongestructureerde prompts gebruiken, terwijl net die promptkwaliteit mee bepaalt of de gegenereerde Python code geschikt is voor data analyse of ETL taken.

Ook recenter onderzoek bevestigt dat promptingtechnieken niet slechts een detail zijn, maar een belangrijk onderdeel van het codegeneratieproces is. In een studie naar codegeneratie vergelijken \textcite{Yuen2025} vier veelgebruikte strategieen, namelijk zero shot, few shot, chain of thought en chain of thought few shot. Zij rapporteren dat strategieen met meer context en expliciete begeleiding beter scoren op functionele correctheid, leesbaarheid en onderhoudbaarheid dan eenvoudige zero shot prompting \autocite{Yuen2025}. Meer bepaald tonen zij aan dat few shot en gecombineerde strategieen vooral nuttig zijn bij complexere programmeertaken, waar het model een voordeel heeft bij voorbeelden of een expliciet redeneerpad. Dit is rechtstreeks toepasbaar op ETL en data transformatie taken, omdat zulke opdrachten meestal uit meerdere opeenvolgende stappen bestaan en daardoor gevoelig zijn voor onvolledige of foutief gestructureerde prompts.

Een bijkomend onderdeel wordt aangehaald door recent onderzoek in \textcite{Yang2025}, waarin prompt engineering wordt beschreven als een methode om niet alleen de functionele juistheid, maar ook de algemene codekwaliteit en ontwikkelervaring te verbeteren \autocite{Yang2025}. Deze studie bespreekt hoe instructie gebaseerde, voorbeeld gebaseerde en hybride prompts elk andere effecten kunnen hebben op de uiteindelijke code, afhankelijk van de complexiteit van de opdracht en het gewenste detailniveau in de output \autocite{Yang2025}. De auteurs benadrukken dat prompt engineering vooral waardevol is wanneer gebruikers een model willen dwingen om explicieter rekening te houden met structuur, randgevallen en verwachtingen rond outputformaat. Voor studenten betekent dit dat prompt engineering niet enkel een techniek is om "meer code" te krijgen, maar een middel om de kans op technisch correcte en bruikbare code te vergroten.

Binnen deze bachelorproef is prompt engineering relevant als theoretisch kader dat verklaart waarom identieke modellen toch sterk verschillende resultaten kunnen opleveren wanneer de prompts anders geformuleerd zijn \autocite{Guo2024,Yuen2025}. De literatuur toont aan dat gestructureerde prompts en prompts met expliciete tussenstappen veelbelovend zijn voor taken waarbij data moet worden ingelezen, opgeschoond, getransformeerd en weggeschreven \autocite{Yuen2025,Yang2025}. Hoewel dit onderzoek zich niet systematisch richt op het vergelijken van verschillende promptstrategieën, blijft prompt engineering een belangrijk theoretisch referentiekader omdat de manier waarop studenten hun vraag formuleren een directe impact heeft op de betrouwbaarheid van de uiteindelijke code.

\subsection{Cybersecurity in LLM gegenereerde code}
\label{sec:cybersecurity_llm_code}

De opkomst van large language models in softwareontwikkeling brengt niet alleen voordelen met zich mee, maar ook belangrijke veiligheidsrisico's. Recent onderzoek toont aan dat LLM gegenereerde code regelmatig kwetsbaarheden bevat, zelfs wanneer het model expliciet gevraagd wordt om veilige code te produceren \autocite{Aldosari2026,Dai2025,Dora2025}. Dit maakt cybersecurity een belangrijk aandachtspunt in onderzoek naar LLM gebaseerde codegeneratie, zeker wanneer studenten gegenereerde code mogelijk overnemen zonder een grondige naar de code te kijken of te testen op veiligheidsrisico's.

Een eerste belangrijk inzicht uit de literatuur is dat veilige codegeneratie niet vanzelf volgt uit algemene codegeneratiecapaciteiten. \textcite{Aldosari2026} onderzoeken hoe prompt engineering kan worden ingezet om kwetsbaarheden in LLM gegenereerde code te verminderen. Hun studie toont aan dat gerichte beveiligingsinstructies in prompts wel degelijk een positief effect kunnen hebben, maar dat dit effect sterk afhangt van de kwaliteit van de prompt en van het type kwetsbaarheid dat wordt onderzocht \autocite{Aldosari2026}. Dit betekent dat een model niet noodzakelijk veiligere code produceert enkel omdat het technisch correcte of goed gestructureerde code genereert. Voor deze bachelorproef is dit relevant, omdat het aantoont dat veiligheid een aparte evaluatiedimensie vormt naast functionele correctheid en performantie.

Een tweede belangrijke bevinding komt uit het onderzoek van \textcite{Dai2025}, dat een reeks recente technieken voor veilige codegeneratie systematisch vergelijkt. De auteurs tonen aan dat veel bestaande methoden die veiligheid proberen te verhogen, dit doen ten koste van de functionele correctheid van de gegenereerde code \autocite{Dai2025}. Wanneer veiligheid en functionaliteit samen op dezelfde gegenereerde code worden geëvalueerd, blijken de verbeteringen vaak beperkt of zelfs misleidend. De studie wijst er bovendien op dat eerdere evaluaties zich te sterk baseerden op één enkele statische analyzer, waardoor echte kwetsbaarheden soms onopgemerkt bleven \autocite{Dai2025}. Dit onderstreept dat een degelijke evaluatie van LLM code zowel functionele tests als meerdere veiligheidscontroles moet combineren.

Ook domeinspecifiek onderzoek naar webapplicatiecode bevestigt dat de verborgen risico's aanzienlijk zijn. \textcite{Dora2025} analyseren door LLM's gegenereerde webapplicatiecode vanuit een security perspectief en stellen vast dat alle onderzochte modellen ernstige zwakheden vertonen op het vlak van authenticatie, sessiebeheer, inputvalidatie en bescherming tegen injectieaanvallen \autocite{Dora2025}. Zelfs wanneer de gegenereerde code op het eerste zicht plausibel en volledig lijkt, blijken fundamentele beveiligingsmechanismen vaak onvolledig of onjuist geïmplementeerd. Hoewel deze studie zich toespitst op webapplicaties, is de onderliggende conclusie ook relevant voor data analyse en ETL contexten: gegenereerde code is niet noodzakelijk veilige code, en onzichtbare fouten kunnen belangrijke gevolgen hebben wanneer scripts toegang krijgen tot databronnen, credentials of productiesystemen.

Een aanvullend argument komt uit recent werk dat de robuustheid van zogenaamde secure code generation methoden verder in vraag stelt. \textcite{Shukla2025} tonen aan dat bestaande technieken voor veilige codegeneratie kwetsbaar blijven onder licht gewijzigde prompts, waardoor eerder gerapporteerde veiligheidswinsten niet altijd standhouden in realistische omstandigheden \autocite{Shukla2025}. Deze bevinding is vooral belangrijk in onderwijscontexten, waar studenten zelden één gestandaardiseerde prompt gebruiken, maar spontaan en variabel formuleren. Daardoor kan de veiligheidskwaliteit van de output niet als stabiel worden verondersteld.

Voor deze bachelorproef betekent dit dat cybersecurity niet als een bijkomstig aspect mag worden behandeld, maar als een volwaardige evaluatiedimensie naast syntaxis, logica en efficiëntie. De studies maken duidelijk dat LLM gegenereerde code technisch correct kan lijken, maar toch kwetsbaar kan zijn door foutieve validatie, onveilige defaults of ontbrekende bescherming tegen misbruik \autocite{Aldosari2026,Shukla2025}. Bij het evalueren van Python code voor data analyse en ETL taken is het daarom zinvol om ook stil te staan bij veiligheidsaspecten zoals omgaan met externe input, veilige bestandsverwerking, correcte behandeling van credentials en het vermijden van gevaarlijke bibliotheek of API patronen.


\subsection{Bestaande evaluatiekaders voor LLM codegeneratie}

Om de prestaties van LLM's op het vlak van codegeneratie te beoordelen, zijn de laatste jaren verschillende benchmarks en evaluatiekaders ontwikkeld \autocite{Ni2024}. Een bekend voorbeeld is het L2CEval framework, dat de taal naar code capaciteiten van LLM's evalueert op basis van een reeks programmeeropdrachten met bijhorende testsuites. \textcite{Ni2024} beschrijven hoe modellen worden beoordeeld op metrics zoals het percentage gegenereerde oplossingen dat alle tests doorstaat en op de gevoeligheid voor varianten in promptformulering \autocite{Ni2024}. Dergelijke benchmarks bieden een gestandaardiseerd zicht op de prestaties van verschillende modellen over een brede variatie aan taken, maar zijn vaak ontworpen met een algemeen doel-groep in gedachten en niet specifiek voor een onderwijscontext of een bepaald toepassingsdomein zoals data engineering.

Naast algemene benchmarks zijn er ook domeinspecifieke evaluaties voorgesteld, bijvoorbeeld voor Data Science taken of agent gebaseerde code generatie in data pipelines \autocite{Nascimento2024}. Deze studies richten zich op het genereren van code voor data manipulatie, modeltraining en evaluatie, en benadrukken dat LLM's niet alleen correct moeten zijn op het niveau van individuele functies, maar ook coherent moeten omgaan met dataflows, dependencies en resources \autocite{Nascimento2024}. Toch blijft het aantal benchmarks dat zich expliciet focust op ETL processen en data engineering in Python beperkt, en wordt er zelden gekeken naar de manier waarop studenten deze tools in de praktijk inzetten of valideren.

\subsection{Python als taal voor data transformatie en ETL}
\label{sec:python-etl}

Een recente systematische literatuurstudie naar Python voor data analyse geeft aan dat Python in zeer veel wetenschappelijke en industriële toepassingen wordt gebruikt voor het volledige traject van gegevensverwerking, van inladen en opschonen tot visualisatie en modellering \autocite{Kabir2024}. De auteurs besluiten dat Python vooral sterk staat door de combinatie van een leesbare syntax, uitgebreide standaardbibliotheken en een rijk ecosysteem aan pakketten voor numerieke berekening en datamanipulatie, waardoor de taal niet alleen geschikt is voor verkennende analyse maar ook voor het opzetten van reproduceerbare data pipelines \autocite{Kabir2024}. Dit maakt Python bijzonder relevant voor omgevingen waar zowel ontwikkelaars als analisten aan dezelfde data pipelines werken.

In een praktijkgericht overzichtswerk over data analyse met Python tonen \textcite{Wei2025} aan hoe Python kan worden ingezet om een volledige workflow op te bouwen waarin data verwerking, visualisatie en voorspelling in dezelfde taal worden beschreven. Zij illustreren hoe typische stappen zoals het inlezen van gegevens uit diverse bronnen, het uitvoeren van transformaties, het samenvoegen van tabellen en het voorbereiden van gegevens voor modellering direct in Python kunnen worden geprogrammeerd, zonder voortdurend tussen verschillende omgevingen te moeten wisselen \autocite{Wei2025}. Deze eenheid in taal en gereedschap verlaagt de kans op fouten bij het overzetten van gegevens en ondersteunt het automatiseren van ETL processen in scripts of notebooks die eenvoudig gedeeld en hergebruikt kunnen worden.

\textcite{Sharma2020} vergelijkt Python expliciet met andere veel gebruikte omgevingen voor data analys. Daaruit blijkt dat Python in het bijzonder voordelen biedt op het vlak van open source beschikbaarheid, community ondersteuning en hoeveelheid gespecialiseerde bibliotheken voor taken zoals datacleaning, feature engineering en model evaluatie \autocite{Sharma2020}. Hoewel andere talen en tools in bepaalde niches efficiënter of eenvoudiger kunnen zijn, concluderen \textcite{Sharma2020} dat Python in de praktijk vaak de voorkeur krijgt omdat het een goede balans biedt tussen gebruiksvriendelijkheid en technische mogelijkheden voor zowel analyse als het bouwen van ETL processen.

Op basis hiervan kan worden gesteld dat Python een onderbouwde keuze is als centrale taal voor data analyse en ETL in deze bachelorproef \autocite{Kabir2024,Dixit2025,Sharma2020}. De taal wordt breed ondersteund in onderzoek en praktijk, laat toe om het volledige pad van ruwe gegevens tot analyse klare dataset in een consistente omgeving uit te werken en sluit nauw aan bij de manier waarop veel moderne data analytics en data engineering projecten worden opgezet.

\subsection{Geheugenverbruik van Python bibliotheken}
\label{sec:geheugenverbruik-python-bibliotheken}

Populaire Python bibliotheken voor data verwerking, zoals \texttt{pandas}, zijn bijzonder gebruiksvriendelijk omdat zij data meestal volledig in het geheugen laden en vervolgens rechtstreeks in tabulaire structuren bewerken \autocite{Breddels2018,Petersohn2020}. Dit model sluit goed aan bij analytische taken die op relatief beperkte datasets worden uitgevoerd, omdat filteren, groeperen, joinen en aggregeren dan zeer expressief en compact kunnen worden beschreven \autocite{Breddels2018}. Tegelijk impliceert deze aanpak dat het volledige datasetvolume in het RAM aanwezig moet zijn, wat bij grotere databestanden snel leidt tot een hogere geheugen gebruik en in sommige gevallen tot performantieproblemen of zelfs het vastlopen van een proces \autocite{Breddels2018,Mekouar2025}.

Onderzoek naar data verwerking in Python toont aan dat het in geheugen werken met tabulaire objecten efficiënt kan zijn zolang de dataset binnen de beschikbare hardwarecapaciteit past, maar dat de schaalbaarheid beperkt wordt wanneer de data groot, breed of complex wordt \autocite{Petersohn2020}. In die situatie ontstaan extra kosten door kopieergedrag, tijdelijke tussenstructuren en het herhaald aanmaken van objecten tijdens transformaties \autocite{Petersohn2020,Mekouar2025}. Vooral bij ETL processen, waar ruwe data vaak eerst moet worden opgeschoond en hervormd voordat ze wordt geladen, kan dit verschil in geheugengebruik een directe impact hebben op de totale uitvoeringstijd en stabiliteit van de pipeline \autocite{Mekouar2025}.

Een bijkomend aandachtspunt is dat de populariteit van \texttt{pandas} soms de indruk wekt dat elke vorm van data verwerking automatisch efficiënt is, terwijl performantie sterk afhangt van de gekozen bewerkingen en de grootte van de dataset \autocite{Breddels2018,Mekouar2025}. Studies die Python gebaseerde dataverwerking vergelijken met alternatieven tonen aan dat het gebruik van in geheugen bibliotheken vooral geschikt is voor interactieve analyse en middelgrote workloads, maar dat voor zeer grote datasets of streaming scenario's andere benaderingen nodig kunnen zijn om geheugenlimieten en bottlenecks te vermijden \autocite{Petersohn2020,Mekouar2025}. In deze bachelorproef is dit relevant omdat ook LLM gegenereerde Python code vaak standaard terugvalt op \texttt{pandas} patronen, terwijl niet elke gegenereerde oplossing rekening houdt met het echte geheugengedrag of met de performantie impact van die keuzes.

Voor ons betekent dit dat de evaluatie van LLM gegenereerde code niet alleen moet nagaan of de output functioneel correct is, maar ook of de voorgestelde data transformaties haalbaar en efficiënt zijn binnen een realistische Python omgeving \autocite{Breddels2018,Petersohn2020,Mekouar2025}. Een model dat een correct antwoord geeft maar tegelijk onnodig veel geheugen gebruikt, kan in een ETL context ongeschikt zijn, zeker wanneer de code later door studenten of in een stageomgeving op grotere datasets wordt toegepast die de code niet meer aankan.


\subsection{Relevantie voor data analyse en ETL in Python}

Data engineering en ETL processen in Python brengen specifieke uitdagingen met zich mee die verder gaan dan het schrijven van losse functies. Typische taken omvatten het inlezen van data uit verschillende bronnen, het uitvoeren van transformaties zoals filteren, join operaties en type conversies, het omgaan met ontbrekende of inconsistente waarden en het laden van de resultaten naar een database of bestandssysteem \autocite{Nascimento2024}. In dit soort workflows zijn niet alleen syntactische correctheid en logica belangrijk, maar ook performantie, schaalbaarheid en robuuste foutafhandeling. Een script dat goed werkt op een klein voorbeeld, kan in de praktijk onbruikbaar blijken wanneer het wordt toegepast op grotere datasets of in een geautomatiseerde pipeline.

De bestaande studies over LLM's en codegeneratie besteedt relatief weinig aandacht aan deze data engineering specifieke aspecten \autocite{Li2023,Ni2024}. Veel benchmarks focussen op algoritmische programmeeropdrachten of relatief kleine codefragmenten, waarbij minder aandacht wordt besteed aan het gebruik van concrete bibliotheken zoals pandas, SQLAlchemy of andere data gerichte tools \autocite{Ni2024}. Ook wordt in de meeste studies niet gekeken naar het gebruik van verouderde of inefficiënte patronen in het Python data ecosysteem, zoals het achterwegen laten van vectorisatie in pandas of het onzorgvuldig omgaan met type conversies. Dit gebrek is juist relevant voor studenten die in hun opleiding leren werken met moderne data engineering praktijken en best practices.

Voor studenten betekent dit dat de bestaande state of the art wel inzicht geeft in algemene sterktes en zwaktes van LLM gegenereerde code, maar weinig concrete handvatten biedt om de kwaliteit van Python code voor data analyse en ETL taken te beoordelen \autocite{Li2023,Esfahani2024}. De stand van zaken in de literatuur toont dus enerzijds aan dat LLM's een krachtig hulpmiddel kunnen zijn, maar anderzijds dat er nog geen duidelijk, praktisch kader bestaat dat hen helpt om de gegenereerde code in hun eigen context systematisch te evalueren.

\subsection{Trainingsdata van coding agents en kwaliteit van de broncode}

Verschillende studies wijzen erop dat de trainingsdata van code georiënteerde LLM’s voornamelijk afkomstig zijn uit grootschalige verzamelingen van open source broncode. Zo beschrijven \textcite{Li2023} dat het model StarCoder werd getraind op The Stack, een dataset van permissively licensed GitHub repositories die meer dan één triljoen tokens aan broncode bevat. Hoewel dergelijke datasets een brede representatie van programmeerstijlen en libraries bieden, bevatten zij onvermijdelijk ook code van wisselende kwaliteit, waaronder experimentele scripts, onvolledige implementaties en voorbeelden uit tutorials. Hierdoor bestaat het risico dat modellen niet alleen goede programmeerpraktijken leren, maar ook minder robuuste patronen reproduceren die in publieke code repositories voorkomen.

Empirisch onderzoek naar LLM gegenereerde code bevestigt deze bezorgdheid. \textcite{Nascimento2024} tonen in hun evaluatie van verschillende LLM gebaseerde codeassistenten voor Data Science taken aan dat gegenereerde oplossingen vaak syntactisch correct zijn, maar niet altijd robuust wanneer ze worden toegepast op complexere datasets of realistische scenario’s. Modellen presteren bijvoorbeeld minder goed wanneer taken grotere datasets, ontbrekende waarden of variabele inputformaten bevatten. Dit suggereert dat de statistische patronen die modellen uit trainingsdata leren niet noodzakelijk leiden tot oplossingen die voldoen aan de kwaliteitsvereisten van productie omgevingen.

Om de prestaties van LLM’s op het vlak van codegeneratie systematisch te beoordelen, zijn de laatste jaren verschillende benchmarks en evaluatiekaders ontwikkeld. \textcite{Ni2024} introduceren bijvoorbeeld SciCode, een benchmark die specifiek gericht is op het evalueren van de programmeercapaciteiten van LLM’s binnen onderzoek en data analysecontexten. Dergelijke benchmarks bouwen voort op eerdere evaluatiesets zoals HumanEval en MBPP en maken gebruik van automatische testcases om te bepalen of gegenereerde code effectief correct functioneert. De resultaten tonen aan dat modellen vaak goed presteren op standaardtaken, maar dat hun prestaties merkbaar dalen wanneer problemen complexer worden of  strengere criteria zoals volledigheid en schaalbaarheid worden toegepast.

Deze bevindingen zijn bijzonder relevant voor domeinen zoals data engineering en ETL processen in Python. In dergelijke contexten volstaat het niet dat code enkel werkt voor een eenvoudig voorbeeld; oplossingen moeten ook schaalbaar, onderhoudbaar en correct functioneren bij grote en heterogene datasets. Zoals \textcite{Nascimento2024} aantonen, genereren LLM’s echter regelmatig oplossingen die eerder lijken op demonstratie of tutorialcode dan op productieklare data pipelines.

In het licht van deze literatuur is het belangrijk te erkennen dat de trainingsdata van code georiënteerde LLM’s slechts in beperkte mate een garantie biedt voor kwalitatieve code. Hoewel deze modellen beschikken over een uitgebreide statistische kennis van programmeertalen en libraries \autocite{Li2023}, ontbreekt een expliciet mechanisme dat formele correctheid, best practices of contextspecifieke kwaliteitsnormen garandeert. Daarom is het noodzakelijk om de kwaliteit van gegenereerde Python code systematisch te evalueren aan de hand van duidelijke testcriteria en kwaliteitsmetingen. De voorliggende bachelorproef sluit hierbij aan door niet alleen te onderzoeken of LLM’s code kunnen genereren voor data analyse en ETL taken, maar ook in welke mate deze code voldoet aan technische kwaliteitscriteria en hoe studenten met behulp van een geschikt evaluatiekader de beperkingen van de onderliggende trainingsdata kunnen compenseren.



\subsection{Positionering van dit onderzoek}

Samenvattend kan worden gesteld dat de huidige stand van zaken drie belangrijke inzichten oplevert. Ten eerste zijn LLM's in staat om snel plausibele Python oplossingen te genereren voor uiteenlopende programmeertaken, maar is de betrouwbaarheid van die oplossingen niet gegarandeerd \autocite{Li2023}. Ten tweede tonen defectanalyses aan dat logische fouten, API misbruik en performantieproblemen frequent voorkomen in LLM gegenereerde code en dat deze fouten niet altijd gemakkelijk te detecteren zijn door minder ervaren programmeurs \autocite{Esfahani2024}. Ten derde bestaan er wel algemene benchmarks en evaluatiekaders, maar ontbreekt een specifiek, onderwijsgericht kader dat zich richt op Python gebaseerde data analyse en ETL taken in een concrete context zoals die van HOGENT studenten \autocite{Ni2024,Nascimento2024}.

Deze bachelorproef positioneert zich precies in deze leegte. Het onderzoek wil de bestaande inzichten uit robuustheidsstudies, defectanalyses en benchmarks vertalen naar een concrete set data engineering taken en evaluatiecriteria die aansluiten bij de leerdoelen en praktijk van studenten Toegepaste Informatica aan HOGENT. Door verschillende LLM's systematisch te evalueren op de kwaliteit van hun Python code voor data analyse en ETL, en de resultaten te vertalen naar een praktisch evaluatiekader en richtlijnen.



\begin{figure}
  \centering
  \includegraphics[width=0.8\textwidth]{grail.jpg}
  \caption[Voorbeeld figuur.]{\label{fig:grail}Voorbeeld van invoegen van een figuur. Zorg altijd voor een uitgebreid bijschrift dat de figuur volledig beschrijft zonder in de tekst te moeten gaan zoeken. Vergeet ook je bronvermelding niet!}
\end{figure}

\begin{listing}
  \begin{minted}{python}
    import pandas as pd
    import seaborn as sns

    penguins = sns.load_dataset('penguins')
    sns.relplot(data=penguins, x="flipper_length_mm", y="bill_length_mm", hue="species")
  \end{minted}
  \caption[Voorbeeld codefragment]{Voorbeeld van het invoegen van een codefragment.}
\end{listing}


\begin{table}
  \centering
  \begin{tabular}{lcr}
    \toprule
    \textbf{Kolom 1} & \textbf{Kolom 2} & \textbf{Kolom 3} \\
    $\alpha$         & $\beta$          & $\gamma$         \\
    \midrule
    A                & 10.230           & a                \\
    B                & 45.678           & b                \\
    C                & 99.987           & c                \\
    \bottomrule
  \end{tabular}
  \caption[Voorbeeld tabel]{\label{tab:example}Voorbeeld van een tabel.}
\end{table}

